
{\actuality} 

\underline{\textbf{Актуальность темы}}

Частота является наиболее точно определяемой физической величиной на сегодняшний день, и при необходимости точного измерения какой-либо физической величины её измерение стараются свести к измерению частоты \cite{Riehle2009}.
Так, с 2019 года все основные величины СИ (кроме секунды) определяются через фундаментальные константы (сейчас они считаются точно известными) и экспериментальное измерение частоты \cite{BIPM2024}.
Секунда, в свою очередь, определяется через частоту радиочастотного перехода между сверхтонкими подуровнями основного состояния в атомах цезия.
Максимально достигнутая относительная неточность определения секунды с помощью цезиевых фонтанов составляет $\sim 10^{-16}$ \cite{petit2015international}. 

Прогресс, однако, не стоит на месте, и за последние десятилетия были созданы реперы частоты, в которых ``маятником'' служат не радиочастотные, а оптические переходы в различных квантовых системах \cite{ludlow2015optical}. 
Увеличение частоты позволяет достигать гораздо лучших значений неточности и нестабильности в определении отношений или разностей частот.
Лучшие из подобных устройств достигают систематической погрешности определения частоты на уровне $10^{-18}$, или даже ниже, а также относительные нестабильности на таком уровне за времена усреднения порядка одного часа \cite{brewer2019al+, bothwell2019jila, bloom2014optical}.
Всё вышеперечисленное позволяет планировать переопределение секунды на основе одного или нескольких оптических часовых переходов, для чего была принята дорожная карта \cite{dimarcq2024roadmap}.
Важным пунктом этой дорожной карты является сравнение различных оптических часов между собой для определения отношения частот часовых переходов с определённой точностью.
Ряд подобных сравнений уже проводился с использованием стабилизированных волоконных линий связи \cite{riehle2017optical, predehl2012920}, однако этот метод не является универсальным и имеет некоторые недостатки.
Например, длина волоконных линий составляла порядка 1000 км, что позволяет сравнивать устройства, разработанные в близкорасположенных институтах, и не позволяет, например, сравнивать оптические часы, расположенные на разных континентах.

Альтернативой этому методу служит построение транспортируемых оптических часов и измерение отношения частот между ними и лабораторными установками.
Подобные устройства также востребованы для решения ряда прикладных задач, таких как задачи релятивистской геодезии \cite{grotti2018geodesy} и спутниковой навигации \cite{schuldt2021optical}.
Разработка транспортируемых оптических часов ведётся рядом лабораторий мира \cite{grotti2018geodesy, cao2022compact, stuhler2021opticlock, semerikov2018compact}, однако создание подобных устройств сопряжено с рядом сложностей.
Так, основной сдвиг частоты часовых переходов в нейтральных атомах, таких как стронций или иттербий, вызывается тепловым излучением окружения \cite{lisdat2021blackbody}. 
Для достижения целевых параметров таких устройств требуется прецизионный контроль температуры окружения или охлаждение устройства до криогенных температур, что приводит к усложнению установки и снижает её надёжность.
Устройства же на основе одиночных ионов требуют гораздо большего времени усреднения для достижения целевых значений неточности и нестабильности.
Всё это служит мотивацией для поиска альтернативных часовых переходов в нейтральных атомах с более низкой чувствительностью частоты перехода к внешним воздействиям.

Одним из таких переходов является часовой переход в нейтральных атомах тулия на длине волны $\lambda = 1.14~\unit{мкм}$ со спектральной шириной $\gamma = 1.2\unit{Гц}$ \cite{golovizin2019inner}.
Этот переход связывает две компоненты тонкой структуры основного состояния атома тулия и меняет состояние $4f$ оболочки, которая находится глубоко внутри заполненных $5s$, $5p$, $6s$ оболочек, из-за чего переход оказывается экранирован от внешних воздействий. 
Интерес к нему возник в 80-х годах XX века, когда было показано, что в атмосфере буферных газов столкновительный сдвиг частоты этого перехода практически отсутствует \cite{александров1983неуширяемая}.
Экспериментальная работа по созданию репера частоты на основе нейтральных атомов тулия ведётся в ФИАНе с 2008 года. 
На текущий момент защищено пять диссертационных работ, посвящённых его разработке, а две, включая эту, находятся в стадии подготовки. 
Диссертация Дениса Сукачёва была посвящена созданию магнито-оптической ловушки для захвата и охлаждения атомов тулия \cite{сукачёв2013лазерное}. 
Вишнякова Гульнара исследовала методы глубокого охлаждения облака атомов с целью понижения их температуры \cite{вишнякова2016вторичное}. 
Работа Артёма Головизина посвящена поиску и прямому возбуждению часового перехода на длине волны $1.14\unit{мкм}$ \cite{головизин2017прямое}.
Фёдорова Елена в своей работе продемонстрировала перезахват атомов в оптическую решётку и их подготовку в определённом квантовом состоянии, которое обеспечивает лучшую стабильность часового перехода \cite{фёдорова2019исследование}. 
Работа Дмитрия Трегубова была посвящена исследованию спектральных характеристик часового перехода в оптической решётке и исследованию основных источников систематических сдвигов \cite{трегубов2020часовой}. 
Находящаяся в стадии подготовки диссертация Дениса Мишина посвящена разработке протокола ``синтетической'' частоты, позволяющего подавить сдвиг, связанный с квадратичным эффектом Зеемана.

Следующим этапом разработки транспортируемых оптических часов на основе атомов тулия является характеризация сдвигов частоты часового перехода, вызываемых оптической решёткой.
В первом приближении, этот сдвиг вызывается электродипольной (E1) поляризуемостью и подавляется при работе на так называемой ``магической'' длине волны оптической решётки (длине волны, на которой электродипольные поляризуемости основного и возбуждённого состояний часового перехода оказываются равными)\cite{katori2003ultrastable, katori2009magic}.
Однако при приближении к уровню неточностей в $10^{-17}$ необходимо учитывать также влияние поляризуемостей старших порядков таких как электроквадрупольная и магнитодипольная поляризуемости, а также гиперполяризуемость. 
Более того, сдвиг, вызываемый ими, зависит не только от глубины оптической решётки и её частоты, но и от движения атомов в оптической решётке \cite{ushijima2018operational}. 
Величина сдвига частоты часового перехода, вызываемого излучением оптической решётки задаётся формулой \ref{eq:LightShift}, в которой $\Delta\nu$ - искомый сдвиг частоты часового перехода, переменными являются глубина решётки и отстройка от магической длины волны, среднее аксиальное колебательное число ансамбля и радиальная температура, а параметрами являются гиперполяризуемость и E2M1 поляризуемость. 
Для характеризации этого сдвига требуется во-первых, уметь подготавливать атомный ансамбля в известном состоянии, а во-вторых, измерение сдвига частоты между двумя ансамблями с разным колебательным числом.

%\begin{equation}
%	\label{eq:LightShift}
%	h\Delta\nu_{LS} (u, \delta_L, n_Z, T_r) \approx (\frac{\partial \alpha^{E1}}{\partial\nu} \delta_L - \alpha^{qm})(n_z + \frac{1}{2})u^{1/2} - [\frac{\partial \alpha^{E1}}{\partial\nu} \delta_L + \frac{3}{2}\beta(n_z^2 + n_z + \frac{1}{2})]u + 2\beta(n_z+\frac{1}{2})u^{3/2} - \beta u^2 , u^j \approx (1 - \frac{j k_B T_r}{uE_r})u^j
%\end{equation}


%На момент начала данной работы использовался следующий протокол подготовки атомов тулия к опросу часового перехода.
%Образец металлического тулия нагревался в атомной печке до температур порядка $T\sim800\unit{К}$ для получения потока атомов. 
%Для их захвата в магнито-оптическую ловушку использовался Зеемановский замедлитель, работающий на длине волны \fst, и позволяет замедлять атомы до скоростей порядка $v\sim10\unit{м/с}.$.
%Первая стадия охлаждения в МОЛ также работает на длине волны \fst, и позволяет охлаждать атомы до температур порядка $T\sim 1\unit{мК}$.
%Это слишком много, поэтому используется вторая стадия охлаждения а 

Используемый на момент начала работы протокол подготовки атомов тулия в начальных состояниях часовых переходов обладал рядом недостатков, которые препятствовали измерению поляризуемостей высоких порядков с достаточной точностью. 
Первоначальный захват и охлаждения атомного ансамбля проводилось с помощью двустадийной магнито-оптической ловушки, что позволяло получать ансамбли с числом атомов порядка $N\sim10^6$ и температурой $T\sim20\unit{мкК}$.
После охлаждения в МОЛ атомы перезахватывались в оптическую решётку на длине волны $\lambda=1063.5\unit{нм}$ с эффективностью перезахвата $\eta\sim1/3$.
Для подготовки атомов на центральном магнитном подуровне использовался метод оптической накачки линейно поляризованным излучением, который приводил к нагреву атомов и потерям из оптической решётки. 
Дополнительное охлаждение проводилось методом ``просеивания'' - кратковременного уменьшения глубины оптической решётки для удаления наиболее горячих атомов из области захвата. 

Данный протокол обладал рядом недостатков:
\begin{itemize}
	\item Оптическая накачка работает на основе рассения фотонов, что приводит к нагреву атомов, и как следствие, к их потерям
	\item Метод охлаждения атомов с помощью просеивания в оптической решётке основан на удалении наиболее горячих атомов из решётки, что, опять же, приводит к их потерям
	\item Метод просеивания не позволял достаточно надёжно охарактеризовать температуру атомов и их распределение по колебательным подуровням оптической решётки
	\item Метод просеивания не позволял достаточно надёжно подготовить атомы в различных колебательных состояниях.
\end{itemize}
Всё это приводит к необходимости разработать новый протокол подготовки атомов тулия в начальных состояниях часового перехода.
%
%Во-первых, оптическая накачка приводила к потерям атомов и их нагреву. 
%Во-вторых, сама суть метода просеивания основана на потерях атомов, что довольно сильно ограничивало число частиц на момент опроса. 
%В-третьих, этот метод не позволял достаточно надёжно охарактеризовать температуру атомов и их распределение по колебательным подуровням оптической решётки.
%В-четвёртых, не получалось достаточно надёжно подготовить атомы в различных колебательных состояниях.
%Всё это приводит к необходимости разработать новый протокол подготовки атомов тулия в начальных состояниях часового перехода.

Есть несколько способов улучшить этот протокол. 
Во-первых, большие потери атомов во время подготовки начальных состтояний можно компенсировать увеличением доли перезахватываемых в решёту атомов.
Добиться этого можно охлаждением ансамбля атомов в МОЛ до меньших температур.
Во-вторых, можно модифицировать протокол подготовки начальных состояний включив в него охлаждение до основного колебательного состояния. 
Этому и посвящена данная диссертация.

%Есть несколько методов охлаждения до основного колебательного состояния.
%В связи с тем, что расстояние между колебательными уровнями ионов в радиочастотных ловушках и атомов в оптических решетках (вдоль оси решетки) составляет порядка $100\unit{кГц}$, наиболее широко применяемым методом глубокого лазерного охлаждения является метод охлаждения на колебательных частотах (sideband cooling) \cite{neuhauser1978optical, perrin1998sideband}. 
%В его основе лежит возбуждение атомного перехода с уменьшением колебательного квантового числа на единицу (красного сайдбенда) и последующий возврат атома в начальное состояние посредством спонтанного распада без изменения колебательного числа (вследствие работы в режиме Лэмба-Дике). 
%Можно выделить три схемы охлаждения на колебательных частотах: (1) с использованием перехода с естественной шириной $\sim 1\div10$\,кГц, которая обеспечивает достаточно быстрый спонтанный распад \cite{ushijima2018operational};
%(2) рамановский переход с использованием различных уровней сверхтонкой структуры либо магнитных подуровней одного уровня \cite{chen2017sympathetic, seck2016raman}; 
%(3) возбуждение узкого часового перехода на красном колебательном сайдбенде с последующей перекачкой населенности \cite{zhang2022subrecoil}. 
%В схемах (2) и (3) для возвращения атома в начальное состояние используется перекачивающее излучение, возбуждающее атом на короткоживущий уровень.
%
%\begin{figure}[t!]
%=	\includegraphics[width=\linewidth]{Chapter0/ОписаниеОхладженияСайдбэндов.pdf}
%	\caption{
%		Различные схемы лазерного охлаждения в оптической решётке до основного колебательного состояния.
%		а) Возбуждение красного колебательного сайдбэнда часового перехода с последующей перекачкой атомов в основное состояние через промежуточный короткоживущий уровень. 
%		Перекачка необходима так как время жизни возбуждённого состояния слишком велико для эффективного охлаждения.
%		б) Охлаждение на спектрально-узком переходе, чья спектральная ширина с одной стороны позволяет селективное возбуждение красного колебательного сайдбэнда (красная линия), и с другой стороны, обеспечивает быстрый спонтанный распад без изменения колебательного числа (черная линия).
%		в) Рамановский переход между двумя (магнитными) подуровнями основного состояния с уменьшением колебательного числа на единицу.
%		г) Характерный спектр перехода в режиме разрешённых боковых полос для двух разных температур атомного ансамбля.
%		Пунктирной линией показан спектр для случая $T_z = T_r = 10\unit{мкК}$ - характерной температуры атомного ансамбля после перезахвата атомов из МОЛ в оптическую решётку.
%		Число атомов в основном колебательном состоянии определяет относительные интенсивности синего и красного сайдбэндов.
%		Сплошная линия соответствует доохлаждённому ансамблю атомов с $T_z = T_r = 1\unit{мкК}$.
%		Большая часть атомов в этом случае расположена в основном колебательном состоянии, что приводит к почти полному исчезновению красной боковой полосы. 
%	} 
%	\label{fig:МетодыОхлаждения}
%\end{figure}
%
%
%\begin{figure}[ht]
%	\centerfloat{
%		\includegraphics[scale=0.9]{Chapter0/Levels.pdf}
%	}
%	\caption{
%		Переходы в нейтральных атомах тулия, использовавшиеся в этой работе.
%	}\label{fig:Levels}
%\end{figure} 


{\aim} данной работы являлось охлаждение атомов тулия в магнито-оптической ловушке на переходе с длиной волны $506.2\unit{нм}$, а также подготовка атомов тулия в основном и первых возбуждённых колебательных состояниях оптической решётки на длине волны $1063.5\unit{нм}$.

Для~достижения поставленной цели необходимо было решить следующие {\tasks}:
\begin{enumerate}[beginpenalty=10000] % https://tex.stackexchange.com/a/476052/104425
  \item Захватить атомы тулия в магнито-оптическую ловушку на переходе с длиной волны $506.2\unit{нм}$, и получить ансамбли атомов тулия с температурой менее $1\unit{мкК}$.
  \item Охарактеризовать влияние оптической решётки на длине волны $1063.5\unit{нм}$ на частоты переходов между различными Зеемановскими компонентами основного и возбуждённого состояния перехода с длиной волны $506.2\unit{нм}$ в атомах тулия.
  \item Охладить атомы тулия, захваченные в оптическую решётку на длине волны $1063.5\unit{нм}$ до основного колебательного состояния оптической решётки. 
  \item Осуществить подготовку атомов тулия в начальных состояниях часовых переходов на колебательных подуровнях $n_z=1$ и $n_z=2$.
\end{enumerate}


{\novelty}
\begin{enumerate}
	\item Продемонстрирован захват атомов тулия в магнито-оптическую ловушку на длине волны 506.2 нм и их охлаждение до температуры $T\sim400~\text{нК}$, а также получены ансамбли атомов тулия с фазовой плотностью $PSD = 3.5\times10^{-4}$ при числе атомов $N = 1.6\times10^6$ и температуре порядка $T\sim\unit{мкК}$.
	\item Измерены динамические скалярная и тензорная поляризуемости  на длине волны $1063.5\unit{нм}$ возбуждённого состояния перехода на длине волны $506.2\unit{нм}$, которые равны соотетственно $\alpha_s = 190(30)~\text{ат.ед}$ и $\alpha_t = 47(7)~\text{ат.ед}$.
	\item Получены ансамбли атомов тулия, захваченные в оптическую решётку с длиной волны $1063.5\unit{нм}$, со средним колебательным числом $n_z = 0.01(1)$.
	\item Получены ансамбли атомов тулия в возбуждённых колебательных состояниях оптической решётки с $n_z = 1$ и $n_z = 2$. Населённость конечного колебательного состояния составляла более $90\%$.
\end{enumerate}

{\influence} 
Разработанный в диссертации протокол подготовки атомов тулия позволит измерить поляризуемости часового перехода высоких порядков с точностью, достаточной для достижения точности разрабатываемого репера частоты на уровне $5\times10^{-18}$, что сравнимо с существующими транспортируемыми оптическими часами. 

%{\methods}

%Для получения представленных результатов производились исследования атомов тулия в магнито-оптической ловушке и в оптической решётке. Экспериментальные данные аппроксимировались теоретическими моделями для определения необходимых характеристик.

{\defpositions}

\begin{itemize}
%	\item Сверхтонкое расщепление уровня $ \ket{4f^{13}(^2F_{7/2}^\text{o})6s6p(^3P_2^\text{o})}$ в нейтральных атомах тулия равняется $\Delta\nu_{hfs} = -3474(2) \unit{МГц}$ %\cite{provorchenko2021investigation}.
	\item При охлаждении нейтральных атомов тулия в магнито-оптической ловушке на переходе $\ket{4f^{13}(^2F^\text{o})6s^2, F = 4} \xrightarrow{} \ket{4f^{13}(^2F_{7/2}^\text{o})6s6p(^3P_2^\text{o}, F=5)}$ с длиной волны $\lambda = 506.2~\unit{нм}$ достижимы температуры облака атомов $T = 400\unit{нК}$, близкие к пределу отдачи используемого перехода ($T_r = 220\unit{нК}$), а также фазовая плотность $\text{PSD}\approx3.5\times10^{-4}$ при числе атомов в МОЛ порядка $N \approx 2\times 10^6$ и температуре $T\approx1\unit{мкК}$ %\cite{provorchenko2023deep}.
	\item Динамические скалярная и тензорная поляризуемости на длине волны $\lambda = 1063.5\unit{нм}$ состояния $\ket{4f^{13}(^2F_{7/2}^\text{o})6s6p(^3P_2^\text{o}), F=5}$ равны $\alpha_s = 190(30)\unit{ат.ед.}$ и $\alpha_t = 47(7)\unit{ат.ед.}$
	\item В случае непараллельных электрического поля оптической решётки и задающего ось квантования магнитного поля, для атомов тулия, удерживающихся в потенциале оптической решётки на длине волны $1063.5\unit{нм}$, возникает эффект конкуренции полей за ось квантования, приводящий к перемешиванию состояний с различными проекциями полного момент для состояния $\ket{4f^{13}(^2F_{7/2}^\text{o})6s6p(^3P_2^\text{o}, F=5, m_{F}=-5)}$ на уровне до 1\%, а также сдвигу часового перехода до $1\unit{Гц}$ для магнитного поля $B = 100\unit{мГс}$ и глубины оптической решётки $U = 1000 \Erec$.
	\item Охлаждение атомов тулия на боковых колебательных частотах перехода \\$\ket{4f^{13}(^2F^\text{o})6s^2, F = 4, m_F=-4} \xrightarrow{} \ket{4f^{13}(^2F_{7/2}^\text{o})6s6p(^3P_2^\text{o}, F=5, m_{F}=-5)}$ с длиной волны $\lambda = 506.2~\unit{нм}$ в оптической решётке на длине волны $\lambda_{lat} = 1063.5\unit{нм}$ позволяет достичь среднего колебательного числа $\bar n_z \approx 0.01(1)$ при числе частиц $N\approx 10^4$ %\cite{provorchenko2024laser}.
	\item Возбуждение ``синих'' боковых частот часового перехода в атомах тулия ($\ket{4f^{13}(^2F^\text{o})6s^2, J=7/2, F = 4(3), m_F = 0} \xrightarrow{} \ket{4f^{13}(^2F^\text{o})6s^2, J=5/2, F=3(2), m_{F}=0}$) позволяет осуществить перенос населённости с основного колебательного состояния на колебательные подуровни $n_z = 1, 2$ с эффективностью 40 и 24 \%, и чистотой конечного состояния $\eta_{n_z=1} = 93(2)\%, \eta_{n_z=2} =95(2)\%$ соответственно %\cite{babichev2025preparation}.
\end{itemize}


{\reliability} полученных результатов подтверждается повторяемостью экспериментальных данных, а также согласованием результатов экспериментов с теоретическими предсказаниями. Результаты находятся в соответствии с результатами, полученными другими авторами.

{\probation}
Основные результаты работы докладывались~на:
\begin{itemize}
	\item D. Provorchenko, A. Golovizin, D Tregubov, D Mishin and N Kolachevsky – “Towards third-stage cooling of thulium atoms with a narrow transition at 506.2 nm” - 30th ANNUAL INTERNATIONAL LASER PHYSICS WORKSHOP – Lyon, France
	\item Д.И. Проворченко, Д.А. Мишин, М.О. Яушев, Д.О. Трегубов, А.А. Головизин, Н.Н. Колачевский – “Лазерное охлаждение атомов тулия на спектрально узком переходе на длине волны 506.2 нм” – “Физика ультрахолодных атомов” - Новосибирск, Россия
	\item D.I. Provorchenko, D.A. Mishin, M.O Yaushev, T.D. Ruzaikin, D.O. Tregubov, A.A. Golovizin, and N.N. Kolachevskiy  - “Narrow-Line Cooling and Initial State Preparation of Thulium Atoms Using the Transition at 506 nm” -- 31th ANNUAL INTERNATIONAL LASER PHYSICS WORKSHOP – Lyon, France
	\item Д.И. Проворченко, Д.А. Мишин, Д.О. Трегубов,
	А.А. Головизин, Н.Н. Колачевский – “Подготовка атомов тулия в начальных состояниях часовых переходов в основном колебательном состоянии оптической решётки с использованием перехода на длине волны 506.2 нм” –  – ЛаПлаз - 2024, МИФИ, - Москва, Россия
\end{itemize}

{\contribution} Все изложенные результаты получены автором лично, либо при его непосредственном участии. 

%
%\ifnumequal{\value{bibliosel}}{0}
%{%%% Встроенная реализация с загрузкой файла через движок bibtex8. (При желании, внутри можно использовать обычные ссылки, наподобие `\cite{vakbib1,vakbib2}`).
%    {\publications} Основные результаты по теме диссертации изложены
%    в~XX~печатных изданиях,
%    X из которых изданы в журналах, рекомендованных ВАК,
%    X "--- в тезисах докладов.
%}%
%{%%% Реализация пакетом biblatex через движок biber
%    \begin{refsection}[bl-author, bl-registered]
%        % Это refsection=1.
%        % Процитированные здесь работы:
%        %  * подсчитываются, для автоматического составления фразы "Основные результаты ..."
%        %  * попадают в авторскую библиографию, при usefootcite==0 и стиле `\insertbiblioauthor` или `\insertbiblioauthorgrouped`
%        %  * нумеруются там в зависимости от порядка команд `\printbibliography` в этом разделе.
%        %  * при использовании `\insertbiblioauthorgrouped`, порядок команд `\printbibliography` в нём должен быть тем же (см. biblio/biblatex.tex)
%        %
%        % Невидимый библиографический список для подсчёта количества публикаций:
%        \printbibliography[heading=nobibheading, section=1, env=countauthorvak,          keyword=biblioauthorvak]%
%        \printbibliography[heading=nobibheading, section=1, env=countauthorwos,          keyword=biblioauthorwos]%
%        \printbibliography[heading=nobibheading, section=1, env=countauthorscopus,       keyword=biblioauthorscopus]%
%        \printbibliography[heading=nobibheading, section=1, env=countauthorconf,         keyword=biblioauthorconf]%
%        \printbibliography[heading=nobibheading, section=1, env=countauthorother,        keyword=biblioauthorother]%
%        \printbibliography[heading=nobibheading, section=1, env=countregistered,         keyword=biblioregistered]%
%        \printbibliography[heading=nobibheading, section=1, env=countauthorpatent,       keyword=biblioauthorpatent]%
%        \printbibliography[heading=nobibheading, section=1, env=countauthorprogram,      keyword=biblioauthorprogram]%
%        \printbibliography[heading=nobibheading, section=1, env=countauthor,             keyword=biblioauthor]%
%        \printbibliography[heading=nobibheading, section=1, env=countauthorvakscopuswos, filter=vakscopuswos]%
%        \printbibliography[heading=nobibheading, section=1, env=countauthorscopuswos,    filter=scopuswos]%
%        %
%        \nocite{*}%
%        %
{\publications} Основные результаты по теме диссертации изложены в~\arabic{citeauthor}~печатных изданиях,
\arabic{citeauthorvak} из которых изданы в журналах, рекомендованных ВАК\sloppy%
%
%\cite{provorchenko2021investigation}
%\cite{mishin2022effect}
%\cite{provorchenko2023deep}
%\cite{provorchenko2024laser}
%\cite{babichev2025preparation}
%        \ifnum \value{citeauthorscopuswos}>0%
%            , \arabic{citeauthorscopuswos} "--- в~периодических научных журналах, индексируемых Web of~Science и Scopus\sloppy%
%        \fi%
%        \ifnum \value{citeauthorconf}>0%
%            , \arabic{citeauthorconf} "--- в~тезисах докладов.
%        \else%
%            .
%        \fi%
%        \ifnum \value{citeregistered}=1%
%            \ifnum \value{citeauthorpatent}=1%
%                Зарегистрирован \arabic{citeauthorpatent} патент.
%            \fi%
%            \ifnum \value{citeauthorprogram}=1%
%                Зарегистрирована \arabic{citeauthorprogram} программа для ЭВМ.
%            \fi%
%        \fi%
%        \ifnum \value{citeregistered}>1%
%            Зарегистрированы\ %
%            \ifnum \value{citeauthorpatent}>0%
%            \formbytotal{citeauthorpatent}{патент}{}{а}{}\sloppy%
%            \ifnum \value{citeauthorprogram}=0 . \else \ и~\fi%
%            \fi%
%            \ifnum \value{citeauthorprogram}>0%
%            \formbytotal{citeauthorprogram}{программ}{а}{ы}{} для ЭВМ.
%            \fi%
%        \fi%
%        % К публикациям, в которых излагаются основные научные результаты диссертации на соискание учёной
%        % степени, в рецензируемых изданиях приравниваются патенты на изобретения, патенты (свидетельства) на
%        % полезную модель, патенты на промышленный образец, патенты на селекционные достижения, свидетельства
%        % на программу для электронных вычислительных машин, базу данных, топологию интегральных микросхем,
%        % зарегистрированные в установленном порядке.(в ред. Постановления Правительства РФ от 21.04.2016 N 335)
%    \end{refsection}%
%    \begin{refsection}[bl-author, bl-registered]
%        % Это refsection=2.
%        % Процитированные здесь работы:
%        %  * попадают в авторскую библиографию, при usefootcite==0 и стиле `\insertbiblioauthorimportant`.
%        %  * ни на что не влияют в противном случае
%        \nocite{vakbib2}%vak
%        \nocite{patbib1}%patent
%        \nocite{progbib1}%program
%        \nocite{bib1}%other
%        \nocite{confbib1}%conf
%    \end{refsection}%
%        %
%        % Всё, что вне этих двух refsection, это refsection=0,
%        %  * для диссертации - это нормальные ссылки, попадающие в обычную библиографию
%        %  * для автореферата:
%        %     * при usefootcite==0, ссылка корректно сработает только для источника из `external.bib`. Для своих работ --- напечатает "[0]" (и даже Warning не вылезет).
%        %     * при usefootcite==1, ссылка сработает нормально. В авторской библиографии будут только процитированные в refsection=0 работы.
%}
%
%При использовании пакета \verb!biblatex! будут подсчитаны все работы, добавленные
%в файл \verb!biblio/author.bib!. Для правильного подсчёта работ в~различных
%системах цитирования требуется использовать поля:
%\begin{itemize}
%        \item \texttt{authorvak} если публикация индексирована ВАК,
%        \item \texttt{authorscopus} если публикация индексирована Scopus,
%        \item \texttt{authorwos} если публикация индексирована Web of Science,
%        \item \texttt{authorconf} для докладов конференций,
%        \item \texttt{authorpatent} для патентов,
%        \item \texttt{authorprogram} для зарегистрированных программ для ЭВМ,
%        \item \texttt{authorother} для других публикаций.
%\end{itemize}
%Для подсчёта используются счётчики:
%\begin{itemize}
%        \item \texttt{citeauthorvak} для работ, индексируемых ВАК,
%        \item \texttt{citeauthorscopus} для работ, индексируемых Scopus,
%        \item \texttt{citeauthorwos} для работ, индексируемых Web of Science,
%        \item \texttt{citeauthorvakscopuswos} для работ, индексируемых одной из трёх баз,
%        \item \texttt{citeauthorscopuswos} для работ, индексируемых Scopus или Web of~Science,
%        \item \texttt{citeauthorconf} для докладов на конференциях,
%        \item \texttt{citeauthorother} для остальных работ,
%        \item \texttt{citeauthorpatent} для патентов,
%        \item \texttt{citeauthorprogram} для зарегистрированных программ для ЭВМ,
%        \item \texttt{citeauthor} для суммарного количества работ.
%\end{itemize}
% Счётчик \texttt{citeexternal} используется для подсчёта процитированных публикаций;
% \texttt{citeregistered} "--- для подсчёта суммарного количества патентов и программ для ЭВМ.
%
%Для добавления в список публикаций автора работ, которые не были процитированы в
%автореферате, требуется их~перечислить с использованием команды \verb!\nocite! в
%\verb!Synopsis/content.tex!.
